% part: intuitionistic-logic
% chapter: semantics
% section: topological-semantics

\documentclass[../../../include/open-logic-section]{subfiles}

\begin{document}

\olfileid[zh]{int}{sem}{top}

\olsection{拓扑语义}

为直觉主义逻辑提供语义的另一种方式，是利用拓扑这一数学概念。

\begin{defn}
  设 $X$ 是一个集合。$X$ 上的一个\emph{拓扑}是一个集合 $\Top{O}
  \subseteq \Pow{X}$，它满足下面这些性质。$\Top{O}$ 的!!{element}被
  称为该拓扑的\emph{开集}。集合 $X$ 连同 $\Top{O}$ 一起被称为一个
  \emph{拓扑空间}。
  \begin{enumerate}
  \item 空集与整个空间都是开的：$\emptyset$、$X \in
    \Top{O}$。
  \item 开集对有穷交封闭：若 $U$、$V \in
    \Top{O}$，则 $U \cap V \in \Top{O}$
  \item 开集对任意并封闭：若对所有 $i \in I$ 都有 $U_i \in
    \Top{O}$，则 $\bigcup \Setabs{U_i}{i \in
      I} \in \Top{O}$。
  \end{enumerate}
\end{defn}

如果开集的集合可以从语境中推知，我们也可以把 $X$ 直接记为某个拓扑；
不过注意，只有在 $X$ 被赋予开集之后，才能把它称为一个拓扑。

\begin{defn}
  直觉主义命题逻辑的一个\emph{拓扑模型}是一个三元组 $\mModel{X} =
  \tuple{X, \Top{O}, V}$，其中 $\Top{O}$ 是 $X$ 上的一个拓扑，而 $V$ 是一个函数，
  它给每个命题变元指派 $\Top{O}$ 中的一个开集。

  给定一个拓扑模型~$\mModel{X}$，归纳定义 $\Prop{X}{!A}$
  如下：
  \begin{enumerate}
  \item $\Prop{X}{\lfalse} = \emptyset$
  \item $\Prop{X}{p} = V(p)$
  \item $\Prop{X}{!A \land !B} = \Prop{X}{!A} \cap \Prop{X}{!B}$
  \item $\Prop{X}{!A \lor !B} = \Prop{X}{!A} \cup \Prop{X}{!B}$
  \item $\Prop{X}{!A \lif !B} = \Interior{(X \setminus \Prop{X}{!A}) \cup
    \Prop{X}{!B}}$
  \end{enumerate}
  这里，$\Interior{V}$ 是一个把集合 $V \subseteq X$
  映到它的\emph{内部}的函数，即它所包含的所有开集的并。换言之，
  \[
  \Interior{V} = \bigcup \Setabs{U}{U \subseteq V \text{ 且 } U \in \Top{O}}.
  \]
\end{defn}

注意，任何集合的内部总是开的，因为它是开集的并。因此，$\Prop{X}{!A}$
总是一个开集。

尽管拓扑语义高度抽象，但仍有几种理解它的方式可能为其提供动因。设
$X$ 的!!{element}，或「点」，是语句可以被评估的那些点。$!A$ 在
其中为真的所有点的集合，就是 $!A$ 所表达的命题。并非每个点集都是一个
潜在的命题；只有 $\Top{O}$ 的!!{element}才是潜在的命题。$!A \Entails !B$ 当且仅当 $!B$
在 $!A$ 为真的每个点处都为真，即对所有~$X$ 都有 $\Prop{X}{!A}
\subseteq \Prop{X}{!B}$。荒谬语句~$\lfalse$
永不为真，所以 $\Prop{X}{\lfalse} = \emptyset$。

为使直觉主义有效的那些规律成立，即要使 $!A \Proves
!B$ 当且仅当 $\Prop{X}{!A} \subset \Prop{X}{!B}$，$!B \land !C$、
$!B \lor !C$ 与 $!B \lif !C$ 所表达的命题，与 $!B$ 和~$!C$
所表达的命题必须有怎样的关系？我们要求对任何 $!A$ 都有 $\lfalse
\Proves !A$，这之所以得到满足，是因为对所有 $U$ 都有
$\emptyset \subseteq U$。由于 $!B \land !C \Proves !B$，我们要求
$\Prop{X}{!B \land !C} \subseteq \Prop{X}{!B}$，同样也要求
$\Prop{X}{!B \land !C} \subseteq \Prop{X}{!C}$。满足 $W \subseteq U$
和 $W \subseteq V$ 的最大集合是 $U \cap V$。
反过来，$!B \Proves !B \lor !C$ 且 $!C \Proves !B \lor !C$，
因此我们要求 $\Prop{X}{!B} \subseteq \Prop{X}{!B \lor !C}$ 和
$\Prop{X}{!C} \subseteq \Prop{X}{!B \lor !C}$。满足 $U \subseteq W$
和 $V \subseteq W$ 的最小集合~$W$ 是 $U \cup V$。

关于 $\lif$ 的定义比较棘手：$!A \lif !B$ 表达的是
与 $!A$ 相结合后衍推 $!B$ 的最弱命题。$!A
\lif !B$ 与 $!A$ 相结合衍推~$!B$ 这一事实，从 $(!A \lif !B)
\land !A \Proves !B $ 可以清楚地看出。因此 $\Prop{X}{!A \lif !B}$
应是满足 $\Prop{X}{!A \lif !B} \cap \Prop{X}{!A}
\subset \Prop{X}{!B}$ 的最大开集，这就引出了我们的定义。

\end{document}
